% !PRINTED 165
\weberpage{165}
%<<CONT
und $C'$. Daraus schliessen wir, dass auch $C$ zu den Invarianten
der Gruppe $S$ gehört, und sprechen dies als Satz aus:

4. \gesperrt{Bildet man aus einer invarianten Form der
Gruppe $S$ beliebige Covarianten, so erhält man
neue invariante Formen der Gruppe.}

\begin{center}
\S.\ 41.

\gesperrt{Der Satz von Hilbert.}
\end{center}

Der Beweis des Satzes von der Endlichkeit des Invariantensystems
 einer linearen Substitutionsgruppe beruht auf einem sehr
allgemeinen Satze über Formensysteme irgend welcher Art, den
\gesperrt{Hilbert} entdeckt und in mannigfachen Untersuchungen über
die Endlichkeit von Invariantensystemen mit ausgezeichnetem
Erfolge angewandt hat, zu dessen Ableitung wir jetzt übergehen
wollen $^{1}$).

I. \gesperrt{Bedeutet $\mathfrak{S}$ irgend ein System von Formen der
$n$ Ver\"anderlichen $x_1,x_2,\ldots x_n$ in endlicher oder
unendlicher Anzahl, so l\"asst sich aus $\mathfrak{S}$ eine
endliche Anzahl von Formen $F_1,F_2,\ldots,F_\mu$ so
ausw\"ahlen, dass jede Form $F$ von $\mathfrak{S}$ durch einen
Ausdruck}
\weq{1}{F = A_1F_1 + A_2F_2 + \cdots + A_\mu F_\mu}
\gesperrt{dargestellt werden kann, worin $A_1,A_2,\ldots A_\mu$
Formen der Variablen $x_1,x_2,\ldots x_n$ sind.}

Die Definition des Formensystemes $\mathfrak{S}$ muss so vollst\"andig
sein, dass von jeder einzelnen Form der Variablen $x$ entschieden
ist, ob sie zu $\mathfrak{S}$ gehört oder nicht, ist aber übrigens an keine
Voraussetzung gebunden.

Besteht $\mathfrak{S}$ nur aus einer endlichen Zahl von Formen, so ist
unser Satz selbstverst\"andlich, denn man kann ja in diesem Falle
die s\"ammtlichen Formen von $\mathfrak{S}$ f\"ur $F_1,F_2,\ldots F_\mu$ nehmen. Der
Beweis wird sich also nur noch mit dem Falle eines unendlichen
Systemes $\mathfrak{S}$ zu befassen haben.

Zur Vereinfachung des Ausdruckes wollen wir das Formensystem
 $F_1,F_2,\ldots F_\mu$ eine \gesperrt{Basis des Systemes} $\mathfrak{S}$ nennen.
Die Functionen $A_1,A_2,\ldots A_\mu$ müssen so beschaffen sein, dass

\bigskip
\noindent\rule{4cm}{0.4pt}

\noindent $^{1}$) \gesperrt{Hilbert}, \glqq Ueber die Theorie der algebraischen Formen\grqq{}. Mathematische
 Annalen, Bd. 36 (1890).
%CONT>>
